\section{Chương~\ref{sec:dofa}. Học với \nt{DOPA}}

Các cơ chế của khớp thần kinh (gói, bộ phát hiện trùng khớp, calci, cửa sổ dẻo) và việc dopamine hành xử như sai số dự đoán đã được đo trực tiếp; việc các quy tắc dẫn tới gradient và cái giá của quảng bá là định lý; kết quả học lựa chọn là tính toán theo mô hình của sách; sơ đồ tính sai số là giả thuyết.

{\footnotesize
\setlength{\tabcolsep}{3pt}\renewcommand{\arraystretch}{1.15}
\begin{longtable}{>{\raggedright}p{0.5cm}>{\raggedright}p{4.0cm}>{\raggedright}p{5.6cm}>{\raggedright}p{2.3cm}>{\raggedright\arraybackslash}p{1.7cm}}
\toprule
TT & Khẳng định & Thí nghiệm và điều đã đo & Nguồn & Trạng thái \\
\midrule\endfirsthead
\toprule
TT & Khẳng định & Thí nghiệm và điều đã đo & Nguồn & Trạng thái \\
\midrule\endhead
1 & Trong não chưa tìm thấy lan truyền ngược sai số ở dạng như trong mạng nhân tạo & Không có thí nghiệm trực tiếp: đây là một khẳng định về sự vắng mặt. Các bài tổng quan phân tích thuật toán đòi hỏi gì (liên kết ngược lặp lại liên kết thuận, một pha riêng) và những xấp xỉ sinh học nào đã được đề xuất. & \cite{Lillicrap2020, WhittingtonBogacz2019} & giả thuyết \\
2 & Trọng số tách thành số vị trí phóng, xác suất phóng và đáp ứng với một gói: $w=N_s\,p\,A_1$~\eqref{eq:quantal} & Khớp thần kinh thần kinh--cơ của ếch khi lượng phóng bị giảm: đáp ứng của bên nhận dao động theo từng bậc, là bội số của đáp ứng thu nhỏ tự phát với một gói; số gói trên một xung tuân theo luật Poisson. & \cite{delCastillo1954} & đã đo \\
3 & Thụ thể của bên nhận là bộ phát hiện trùng khớp; calci khởi động sự thay đổi trọng số & Kẹp miếng màng (patch-clamp) nơ-ron chuột nhắt: ion Mg$^{2+}$ bịt kênh do glutamate mở ở điện thế nghỉ và rời đi khi khử cực. Lát cắt hồi hải mã: chất tạo chelat calci trong bên nhận chặn tăng cường dài hạn, còn giải phóng calci bằng ánh sáng bên trong tế bào tự nó làm tăng truyền dẫn. & \cite{Nowak1984, Malenka1988calcium} & đã đo \\
4 & Phía nào cũng có thể thực thi quyết định: $A_1$ tăng ở bên nhận, $p$ đổi ở bên gửi & Glutamate được giải phóng bằng ánh sáng phía trên một gai gây ra sự tăng trưởng nhanh và bền của gai cùng dòng qua các thụ thể của nó; sự tăng cường đi kèm việc gắn thêm thụ thể AMPA. Khử cực bên nhận tạm thời ức chế sự phóng ở bên gửi; hiệu ứng được truyền bởi endocannabinoid đi ngược qua khe (đã chỉ ra ở các đầu vào \nt{GABA}). Suy giảm ngắn hạn phụ thuộc vào xác suất phóng ở bên gửi. & \cite{Matsuzaki2004, Malinow2002, WilsonNicoll2001, Tsodyks1997} & đã đo \\
5 & Khớp thần kinh mạnh lên khi bên gửi và bên nhận cùng hoạt động~\eqref{eq:hebb} & Thỏ gây mê: sau những loạt xung ngắn theo đường vào của hồi hải mã, đáp ứng được tăng cường trong khoảng từ $30$~phút tới $10$~giờ. Trong lát cắt hồi hải mã, xung của bên nhận chỉ làm mạnh những khớp thần kinh đang hoạt động cùng lúc. & \cite{Bliss1973LTP, Kelso1986hebb} & đã đo \\
6 & Quy tắc~\eqref{eq:oja} tìm vectơ riêng chính của tương quan đầu vào (mệnh đề~\ref{prop:oja}) & Định lý; chứng minh trong chương. Chưa có thí nghiệm nào cho thấy một nơ-ron học được thành phần chính của các đầu vào của nó; cơ chế giới hạn sự tăng trọng số trong khớp thần kinh chưa được xác định. & \cite{Oja1982} & lý thuyết \\
7 & Hebb theo thời gian: cửa sổ khoảng $\pm20\ms$ (hình~\ref{fig:stdp}); từ các chuỗi lặp lại mọc ra các chuỗi liên kết & Các cặp nơ-ron hồi hải mã trong mẫu nuôi cấy: xung của bên nhận trong vòng $20\ms$ sau xung của bên gửi cho tăng cường bền, trong vòng $20\ms$ trước thì cho suy yếu; cả hai phụ thuộc thụ thể NMDA. Tỷ lệ $A_-=0.6A_+$ trên hình chỉ để minh họa. Việc các chuỗi mọc ra trong mạng theo cách này chưa được đo trực tiếp. & \cite{BiPoo1998} & đã đo \\
8 & Vết~\eqref{eq:threefactor}: dopamine đến sau hoạt động đồng thời $1$--$2\,\text{s}$ thì làm mạnh liên kết, muộn hơn thì không (mệnh đề~\ref{prop:threefactor}) & Lát cắt thể vân, kích thích quang di truyền riêng rẽ đầu vào glutamate và dopamine: gai chỉ lớn lên nếu dopamine đến sau đầu vào glutamate $0.3$--$2\,\text{s}$; cửa sổ do sự tăng nhanh và phân rã của cAMP quyết định. Ở vỏ não, việc ghép cặp để lại những vết riêng cho tăng cường và suy yếu, và các thụ thể monoamine biến chúng thành thay đổi dài hạn trong cửa sổ cỡ giây. & \cite{Yagishita2014, He2015eligibility} & đã đo \\
9 & Cửa sổ dài cỡ giây cũng có khi không có dopamine: tín hiệu thứ ba là sự kích thích mạnh của bên nhận & Chuột nhắt sống, vùng CA1: một sự kiện điện thế bình nguyên ở bên nhận làm mạnh các đầu vào đến trong vài giây trước và sau nó, và chỉ sau một lượt đi qua đã tạo ra một trường vị trí. & \cite{Bittner2017} & đã đo \\
10 & Bước trung bình của quy tắc ba thừa số đi theo gradient phần thưởng~\eqref{eq:perturbation} (mệnh đề~\ref{prop:gradient}) & Định lý gradient cho việc học theo độ lệch ngẫu nhiên; trong chương được suy ra cho nhiễu nhỏ. & \cite{Williams1992} & lý thuyết \\
11 & Nhiễu là sự tìm kiếm: mạng học từ những độ lệch ngẫu nhiên của chính nó & Chim sẻ vằn non: tắt nhân đầu ra của mạch hạch nền làm mất đi đột ngột và có hồi phục tính biến thiên của tiếng hót. Chim sẻ vằn trưởng thành: nhiễu được đưa vào những lần hót một âm tiết có cao độ nằm về một phía so với trung bình, và chim nhanh chóng dịch cao độ sang phía kia --- học từ độ phân tán của chính mình. & \cite{Olveczky2005, TumerBrainard2007} & đã đo \\
12 & Lựa chọn giữa hai phương án được học khi $\tau_e=1\,\text{s}$ và $\delta=R-\bar R$; không học được khi $\tau_e=50\ms$; không trừ kỳ vọng thì trọng số tăng không giới hạn, còn với tốc độ học lớn mạng có thể cố định lựa chọn tệ hơn & \texttt{models/dofa\_learning.py}, $\eta=0.05$, $500$ lượt thử, $6$ seed. $\tau_e=1\,\text{s}$: tỷ lệ chọn $A$ $0.65$--$0.79\to0.96$--$1.00$, trọng số $0.85$--$1.33$. $\tau_e=50\ms$: $0.38$--$0.55$, trọng số không đổi. $\delta=R$: trọng số của bên thắng $860$--$1190$. $\delta=R$, $\eta=0.5$, $3$ seed: một lần bị cố định ở $B$ ($0.04\to0$). Script in ra cả bốn trường hợp; tính lại khớp với chương. & \texttt{models/\allowbreak dofa\_\allowbreak learning.py} & mô hình \\
13 & Thời gian học bằng một tín hiệu chung tăng tỷ lệ với số nơ-ron gây nhiễu (mệnh đề~\ref{prop:broadcastcost}) & Đường cong học của mạng tuyến tính: nhiễu loạn nơ-ron chậm hơn gradient chính xác một số lần bằng số nơ-ron đầu ra, nhiễu loạn trọng số còn chậm thêm một số lần bằng số đầu vào. & \cite{Werfel2005} & lý thuyết \\
14 & Đầu ra của \nt{DOPA} đi trước hết tới các vùng chọn hành động & Tái dựng toàn bộ sợi trục của từng nơ-ron \nt{DOPA} đường liềm đen--thể vân ở chuột cống: các nhánh của một tế bào phủ $0.45$--$5.7\,\%$ thể tích thể vân (trung bình $2.7\,\%$), ngoài thể vân gần như không có nhánh. & \cite{Matsuda2009} & đã đo \\
15 & Vỏ não học cách dự đoán đầu vào của nó, sai số được tính tại chỗ & Tổng quan: trong vỏ não cảm giác có những đáp ứng với sự không khớp giữa đầu vào được kỳ vọng và đầu vào thực đến. Việc đây là quy tắc học chung của vỏ não chưa được chứng minh. & \cite{Keller2018} & giả thuyết \\
16 & Dopamine hành xử như sai số~\eqref{eq:ctd}: đỉnh vọt, dịch chuyển sang tín hiệu báo trước, hõm sụt (hình~\ref{fig:ctdcases}) & Nơ-ron \nt{DOPA} ở khỉ: đỉnh vọt khi có nước ép bất ngờ; sau khi học --- đỉnh vọt lúc có tín hiệu báo trước và không đáp ứng với nước ép đã hứa; hõm sụt khi nước ép không đến. Dạng liên tục~\eqref{eq:ctd} là lý thuyết. & \cite{Schultz1997, Doya2000} & đã đo \\
17 & Sơ đồ tính $\dd V/\dd t$ và trừ đi kỳ vọng & Chuột nhắt, ghi và quang di truyền ở vùng mái bụng: nơ-ron \nt{DOPA} trừ kỳ vọng khỏi phần thưởng; các nơ-ron \nt{GABA} lân cận ức chế chúng khi phần thưởng được kỳ vọng, và kích thích hay ức chế các nơ-ron này làm thay đổi sai số. Đạo hàm được tính thế nào thì chưa được chỉ ra. & \cite{Eshel2015arithmetic} & giả thuyết \\
18 & Nơ-ron \nt{DOPA} là bộ tạo nhịp; hõm sụt nông hơn đỉnh vọt & Trong lát cắt, nơ-ron \nt{DOPA} phát xung tự phát và rất đều, trên nền một điện thế tạo nhịp chậm. Ở khỉ, tần số bám theo sai số dương một cách định lượng, còn sai số âm thì chỉ truyền được yếu ớt. & \cite{GraceOnn1989, Bayer2005} & đã đo \\
19 & Bộ hành động và bộ phê bình (hình~\ref{fig:actorcritic}) & Thuật toán lấy từ lý thuyết học tăng cường. Ở người (fMRI), thể vân bụng hành xử như bộ phê bình cả khi có lựa chọn lẫn khi không, thể vân lưng --- chỉ khi phải chọn hành động. & \cite{SuttonBarto2018, ODoherty2004actor} & giả thuyết \\
20 & Dấu của $\delta$ trong quy tắc bị đảo ở các nơ-ron có họ thụ thể thứ hai & Lát cắt thể vân: ở nơ-ron có thụ thể D1, việc ghép cặp cho tăng cường, cần D1; ở nơ-ron có D2 --- suy yếu, cần D2. & \cite{Shen2008} & đã đo \\
\bottomrule
\end{longtable}}
